\chapter{Introducción}

El análisis de imágenes médicas es un campo donde la ingeniería informática puede aportar herramientas de apoyo para tareas de segmentación, medición, visualización y documentación. En particular, la resonancia magnética (RM) de columna lumbar permite observar estructuras como cuerpos vertebrales, discos intervertebrales y canal espinal, cuyo análisis requiere criterio profesional y una correcta interpretación anatómica. En este contexto existe una oportunidad de asistencia informática: automatizar tareas operativas que hoy pueden requerir revisión manual o semimanual, sin sustituir la decisión del especialista. La segmentación anatómica asistida por inteligencia artificial puede servir como base para visualizar regiones de interés, calcular mediciones geométricas simples y organizar resultados de manera trazable.

El presente Proyecto Final de Ingeniería propone el diseño y desarrollo de un prototipo funcional de software para RM lumbar. Inicialmente centrado en el plano sagital, el proyecto incorporó además el plano axial como complemento, a partir de un hallazgo del relevamiento con profesionales, y admite estudios en formatos estándar de imagen médica (DICOM, .ima o .mha). El MVP contempla la carga y normalización de estudios, la segmentación de vértebras, discos intervertebrales y canal espinal, la extracción de mediciones cuantitativas acotadas, la visualización superpuesta de máscaras y la generación de una salida estructurada editable, incluido un preinforme automático revisable. El sistema se plantea como herramienta de apoyo: no emite diagnósticos, no recomienda tratamientos y no reemplaza al profesional. El desarrollo se apoya principalmente en datasets públicos anotados, especialmente SPIDER para el plano sagital y el conjunto de Sudirman / Al-Kafri para el plano axial, lo que favorece la reproducibilidad académica y permite validar técnicamente las segmentaciones contra máscaras de referencia. Esta decisión también reduce la dependencia inicial de datos clínicos privados y mantiene el alcance dentro de una prueba de concepto verificable.

La organización del documento avanza desde la definición del problema hacia la fundamentación y el diseño del prototipo. Los primeros capítulos presentan el problema, la justificación, los objetivos, el alcance y la solución propuesta. Luego se desarrollan el estado del arte y el marco teórico, que establecen los antecedentes y conceptos necesarios para comprender las decisiones técnicas del proyecto. En etapas posteriores del documento final se incorporarán user research, requerimientos, arquitectura, metodología, implementación, validación, resultados y conclusiones.
